\section{Related Work and Research Gap}
\label{sec:related}
\subsection{AES implementations and workload characteristics}
AES uses a fixed block size and supports 128-, 192- and 256-bit keys, with 10, 12 and 14 rounds, respectively~\cite{fips197}. The mode of operation determines how these block transformations are combined. CBC encryption introduces dependencies between successive blocks, whereas CTR generates keystream from independent counter values~\cite{sp80038a}. GCM combines counter-mode encryption with authentication and includes tag generation or verification~\cite{sp80038d}. These differences affect the available parallelism and the work performed per message, so performance must be interpreted for a specified mode and operation.

Research on efficient AES implementations has examined both instruction support and software organization. K\"asper and Schwabe~\cite{kasper2009} presented a bitsliced implementation for 64-bit Intel processors and investigated AES-GCM performance together with resistance to timing attacks. Their work demonstrates that implementation choices matter for both efficiency and security. Intel's AES-NI instructions provide dedicated operations for AES processing~\cite{intelaesni}. Such mechanisms explain why the same standardized cipher can have substantially different execution costs across implementations, even before application overhead is considered.

Library interfaces add another layer. OpenSSL EVP exposes cipher operations through an abstraction whose selected implementation depends on the library configuration and available processor features~\cite{openssl}. PyCryptodome provides AES modes through a Python interface backed by native code~\cite{pycryptodome}. Neither the common algorithm name nor the availability of processor instructions guarantees that two interfaces select equivalent code or include the same work in a timed call. The present comparison therefore concerns the particular library configurations and interfaces recorded in the experiments.

\subsection{Benchmark interpretation and the scope of this study}
Systems research has established that apparently minor experimental choices can affect measured performance. Mytkowicz et al.~\cite{mytkowicz2009} demonstrated measurement bias arising from experimental setup. Kalibera and Jones~\cite{kalibera2013} developed a methodology that accounts for variation at different levels of repetition. These studies motivate documenting the execution environment and distinguishing repeated timings within a session from replication across sessions or machines.

For encryption benchmarks, measurement scope introduces a further interpretive issue. Timing the cipher transformation answers a question about processing prepared data. Timing a surrounding software operation also captures some combination of preparation, allocation, interface overhead and result handling. In authenticated encryption, tag processing and associated data introduce additional work. A throughput comparison is therefore meaningful only when these inclusions are identified alongside the payload size and implementation.

The practical question addressed here is how strongly implementation and timing choices alter the interpretation of CPU encryption benchmarks. This study connects AES implementation research with benchmarking methodology through documented comparisons of library paths, timing scopes and cryptographic workloads. Its contribution is empirical clarification of these relationships on two systems, rather than a new AES implementation.
